\begin{textsample}{POS Dim 2 – unprofiled\_gpt – Score 9.00 – unprofiled\_gpt/t001509\_gpt.txt}  \label{ex:f2_pos_021}
There are a few fairly basic theories in careers guidance, and they ’re really all about how people make decisions.
One that ’s been very influential, and is used a \textbf{lot} in careers services, is psychologically based.
It assumes that once you ’ve reached a certain level academically, you \textbf{become} \textbf{more} independent and can make choices grounded in your \textbf{own} skills, your values, and what genuinely interests you.
That ’s the ideal model, if you like, and much of our practice is \textbf{built} around helping students clarify those things about \textbf{themselves}.
But you \textbf{also} have to be realistic and recognise that not everyone actually makes decisions on that basis.
Some students will appear to be going through that reflective, individual process, but in \textbf{fact} their choices are driven by quite \textbf{different} factors, and you have to keep that in mind when you ’re working with them.

% matched lemmas: also, become, build, different, fact, lot, more, own, themselves
\end{textsample}
